\documentclass[main.tex]{subfiles}
\begin{document}\newpage
\setdoublesep{0.35700 em}  % 'Bond Spacing'
\setatomsep{1.78500 em}    % 'Fixed Length'
\setbondoffset{0.18265 em} % 'Margin Width'
\newcommand{\bondwidth}{0.06642 em} % 'Line Width'
\setbondstyle{line width = \bondwidth}
\newgeometry{left=0.8in,right=0.8in, top=2.5cm,bottom=2cm}
\fancyhfoffset[E,O]{0pt}
\setlength{\columnsep}{30pt}
\begin{conclusion}
\end{conclusion}
\setstretch{0.3}
\begin{multicols*}{2}

%{\raggedright\textsc{\textbf{Electron-Dot structures of atoms }}\par}
%
%\begin{enumerate}
%
%\item The electron-dot structure of \ce{H^{}} is:
%\begin{enumerate}[label=(\alph*)]
%\begin{multicols*}{2}
%\item \hspace{.05in}\lewis{0.,H}\hspace{.05in}
%\item \hspace{.05in}\lewis{0.4.,H}\hspace{.05in}
%\item \hspace{.05in}\lewis{0.4.6.,H}\hspace{.05in}
%\item \hspace{.05in}\lewis{0.2.4.6.,H}\hspace{.05in}
%\item \hspace{.05in}\lewis{0.2:4.6.,H}\hspace{.05in}
%\end{multicols*}\flushright  {\small Ans: (a)}
%\end{enumerate}
%
%\item The electron-dot structure of \ce{Mg^{}} is:
%\begin{enumerate}[label=(\alph*)]
%\begin{multicols*}{2}
%\item \hspace{.05in}\lewis{0.,Mg}\hspace{.05in}
%\item \hspace{.05in}\lewis{0.4.,Mg}\hspace{.05in}
%\item \hspace{.05in}\lewis{0.4.6.,Mg}\hspace{.05in}
%\item \hspace{.05in}\lewis{0.2.4.6.,Mg}\hspace{.05in}
%\item \hspace{.05in}\lewis{0.2:4.6.,Mg}\hspace{.05in}
%\end{multicols*}\flushright  {\small Ans: (b)}
%\end{enumerate}
%\item The electron-dot structure of \ce{B^{}} is:
%\begin{enumerate}[label=(\alph*)]
%\begin{multicols*}{2}
%\item \hspace{.05in}\lewis{0.,B}\hspace{.05in}
%\item \hspace{.05in}\lewis{0.4.,B}\hspace{.05in}
%\item \hspace{.05in}\lewis{0.4.6.,B}\hspace{.05in}
%\item \hspace{.05in}\lewis{0.2.4.6.,B}\hspace{.05in}
%\item \hspace{.05in}\lewis{0.2:4.6.,B}\hspace{.05in}
%\end{multicols*}\flushright  {\small Ans: (c)}
%\end{enumerate}
%
%%\item The electron-dot structure of \ce{N^{}} is:
%%\begin{enumerate}[label=(\alph*)]
%%\begin{multicols*}{2}
%%\item \hspace{.05in}\lewis{0.,N}\hspace{.05in}
%%\item \hspace{.05in}\lewis{0.4.,N}\hspace{.05in}
%%\item \hspace{.05in}\lewis{0.4.6.,N}\hspace{.05in}
%%\item \hspace{.05in}\lewis{0.2.4.6.,N}\hspace{.05in}
%%\item \hspace{.05in}\lewis{0.2:4.6.,N}\hspace{.05in}
%%\end{multicols*}\flushright  {\small Ans: (e)}
%%\end{enumerate}
%
%\item The electron-dot structure of \ce{O^{}} is:
%\begin{enumerate}[label=(\alph*)]
%\begin{multicols*}{2}
%\item \hspace{.05in}\lewis{0.,O}\hspace{.05in}
%\item \hspace{.05in}\lewis{0.4.,O}\hspace{.05in}
%\item \hspace{.05in}\lewis{0.4.6.,O}\hspace{.05in}
%\item \hspace{.05in}\lewis{0.2.4.6.,O}\hspace{.05in}
%\item \hspace{.05in}\lewis{0.2:4.6:,O}\hspace{.05in}
%\end{multicols*}\flushright  {\small Ans: (e)}
%\end{enumerate}
%
%\item The electron-dot structure of \ce{C^{}} is:
%\begin{enumerate}[label=(\alph*)]
%\begin{multicols*}{2}
%\item \hspace{.05in}\lewis{0.,C}\hspace{.05in}
%\item \hspace{.05in}\lewis{0.4.,C}\hspace{.05in}
%\item \hspace{.05in}\lewis{0.4.6.,C}\hspace{.05in}
%\item \hspace{.05in}\lewis{0.2.4.6.,C}\hspace{.05in}
%\item \hspace{.05in}\lewis{0.2:4.6:,C}\hspace{.05in}
%\end{multicols*}\flushright  {\small Ans: (d)}
%\end{enumerate}
%
%\item The electron-dot structure of \ce{F^{-1}} is:
%\begin{enumerate}[label=(\alph*)]
%\begin{multicols*}{2}
%\item \hspace{.05in}\lewis{0.2.4.6.,F}\hspace{.05in}\ce{^{-}}
%\item \hspace{.05in}\lewis{0:4:,F}\hspace{.05in}\ce{^{-}}
%\item \hspace{.05in}\lewis{0:4:6.,F}\hspace{.05in}\ce{^{-}}
%\item \hspace{.05in}\lewis{0.2.4.6.,F}\hspace{.05in}\ce{^{-}}
%\item \hspace{.05in}\lewis{0:2:4:6:,F}\hspace{.05in}\ce{^{-}}
%\end{multicols*}\flushright  {\small Ans: (e)}
%\end{enumerate}
%
%\item The electron-dot structure of \ce{N^{-3}} is:
%\begin{enumerate}[label=(\alph*)]
%\begin{multicols*}{2}
%\item \hspace{.05in}\lewis{0.2.4.6.,N}\hspace{.05in}\ce{^{-3}}
%\item \hspace{.05in}\lewis{0:4:,N}\hspace{.05in}\ce{^{-3}}
%\item \hspace{.05in}\lewis{0:4:6.,N}\hspace{.05in}\ce{^{-3}}
%\item \hspace{.05in}\lewis{0.2.4.6.,N}\hspace{.05in}\ce{^{-3}}
%\item \hspace{.05in}\lewis{0:2:4:6:,N}\hspace{.05in}\ce{^{-3}}
%\end{multicols*}\flushright  {\small Ans: (e)}
%\end{enumerate}

{\raggedright\textsc{\textbf{Electron-Dot structures of molecules }}\par}
\begin{enumerate}

\item The electron-dot structure of HI is:
\begin{enumerate}[label=(\alph*)]
\begin{multicols*}{2}
\item \chemfig{ \lewis{46,H}-[:0]\lewis{026,I}}\hspace{.05in}
\item \chemfig{ \lewis{46,H}-[:0]\lewis{26,I}}\hspace{.05in}
\item \chemfig{ \lewis{46,H}-[:0]\lewis{6,I}}\hspace{.05in}
  
\item \chemfig{ \lewis{26,H}-[:0]\lewis{26,I}}\hspace{.05in}
\item \chemfig{ H-[:0]\lewis{026,I}}\hspace{.05in}
\end{multicols*}\flushright  {\small Ans: (e)}
\end{enumerate}


\item Draw the electron-dot structure of \ce{BH3}:
\begin{center}\begin{tikzpicture}
\draw[step=.5cm,gray,very thin] (-2,-2) grid (2,2);
\end{tikzpicture}\end{center}

\item Draw the electron-dot structure of \ce{CH4}:
\begin{center}\begin{tikzpicture}
\draw[step=.5cm,gray,very thin] (-2,-2) grid (2,2);
\end{tikzpicture}\end{center}

\item Draw the electron-dot structure of \ce{NH4^{+}}:
\begin{center}\begin{tikzpicture}
\draw[step=.5cm,gray,very thin] (-2,-2) grid (2,2);
\end{tikzpicture}\end{center}

\item Draw the electron-dot structure of \ce{H3O^{+}}:
\begin{center}\begin{tikzpicture}
\draw[step=.5cm,gray,very thin] (-2,-2) grid (2,2);
\end{tikzpicture}\end{center}




\item Indicate the charge of the C atoms in the following electron-dot structure:\\
\begin{center}
\chemleft[
\chemfig{\lewis{2,C}( (-[:0]H)(-[:180]H)(-[:270]H))}\chemright]
\chemmove{\node[] at (0pt,15pt) {\footnotesize $2-$};}
\end{center}
\begin{enumerate}[label=(\alph*)]
\begin{multicols*}{2}
\item +1
\item +2
\item -1
\item -2
\item 5
\end{multicols*}\flushright  {\small Ans: (d)}
\end{enumerate}

\item Indicate the charge of the O atoms in the following electron-dot structure:\\
\begin{center}
\chemleft[
\chemfig{\lewis{2,O}( (-[:0]H)(-[:180]H)(-[:270]H))}\chemright]
\chemmove{\node[] at (0pt,15pt) {\footnotesize $+$};}
\end{center}
\begin{enumerate}[label=(\alph*)]
\begin{multicols*}{2}
\item +1
\item +2
\item -1
\item -2
\item 4
\end{multicols*}\flushright  {\small Ans: (a)}
\end{enumerate}

\item Indicate the charge of the O atoms in the following electron-dot structure:\\
\begin{center}
\chemleft[
\chemfig{N( (-[:0]H)(-[:90]H)(-[:180]H)(-[:270]H))}\chemright]
\chemmove{\node[] at (0pt,35pt) {\footnotesize $+$};}
\end{center}
\begin{enumerate}[label=(\alph*)]
\begin{multicols*}{2}
\item +1
\item +2
\item -1
\item -2
\item 4
\end{multicols*}\flushright  {\small Ans: (a)}
\end{enumerate}

\item Indicate the charge of the O atoms in the following electron-dot structure:\\
\begin{center}
\chemleft[
\chemfig{C((-[:90]H)(-[:180]H)(-[:270]H))-[:0]\lewis{026,O}}\chemright]
\chemmove{\node[] at (0pt,35pt) {\footnotesize $-$};}
\end{center}
\begin{enumerate}[label=(\alph*)]
\begin{multicols*}{2}
\item +1
\item +2
\item -1
\item -2
\item 4
\end{multicols*}\flushright  {\small Ans: (c)}
\end{enumerate}



{\raggedright\textsc{\textbf{Molecular shape}}\par}

\item  Identify the molecular shape of the molecule: \ce{H2}
\begin{enumerate}[label=(\alph*)]
\begin{multicols*}{2}
\item Linera
\item Trigonal planar
\item Bent
\item Tetrahedral
\item  Trigonal pyramidal
\end{multicols*}\flushright  {\small Ans: (a)}
\end{enumerate}

\item  Identify the molecular shape of the molecule: \ce{BeCl2}
\begin{enumerate}[label=(\alph*)]
\begin{multicols*}{2}
\item Linera
\item Trigonal planar
\item Bent
\item Tetrahedral
\item  Trigonal pyramidal
\end{multicols*}\flushright  {\small Ans: (a)}
\end{enumerate}

\item  Identify the molecular shape of the molecule: \ce{BF3}
\begin{enumerate}[label=(\alph*)]
\begin{multicols*}{2}
\item Linera
\item Trigonal planar
\item Bent
\item Tetrahedral
\item  Trigonal pyramidal
\end{multicols*}\flushright  {\small Ans: (b)}
\end{enumerate}

\item  Identify the molecular shape of the molecule: \ce{NH3}
\begin{enumerate}[label=(\alph*)]
\begin{multicols*}{2}
\item Linera
\item Trigonal planar
\item Bent
\item Tetrahedral
\item  Trigonal pyramidal
\end{multicols*}\flushright  {\small Ans: (e)}
\end{enumerate}

\item  Identify the molecular shape of the molecule: \ce{CH4}
\begin{enumerate}[label=(\alph*)]
\begin{multicols*}{2}
\item Linera
\item Trigonal planar
\item Bent
\item Tetrahedral
\item  Trigonal pyramidal
\end{multicols*}\flushright  {\small Ans: (d)}
\end{enumerate}

{\raggedright\textsc{\textbf{Polarity}}\par}

\item  The \ce{H2O} molecule is
\begin{enumerate}[label=(\alph*)]
\begin{multicols*}{2}
\item polar
\item nonpolar
\end{multicols*}\flushright  {\small Ans: (a)}
\end{enumerate}

\item  The \ce{HCl} molecule is
\begin{enumerate}[label=(\alph*)]
\begin{multicols*}{2}
\item polar
\item nonpolar
\end{multicols*}\flushright  {\small Ans: (a)}
\end{enumerate}

\item  The \ce{H2} molecule is
\begin{enumerate}[label=(\alph*)]
\begin{multicols*}{2}
\item polar
\item nonpolar
\end{multicols*}\flushright  {\small Ans: (b)}
\end{enumerate}

\item  The \ce{NH3} molecule is
\begin{enumerate}[label=(\alph*)]
\begin{multicols*}{2}
\item polar
\item nonpolar
\end{multicols*}\flushright  {\small Ans: (a)}
\end{enumerate}

{\raggedright\textsc{\textbf{Hybrid Orbitals}}\par}

\item  Indicate the hybridization of: \ce{NH3} 
\begin{enumerate}[label=(\alph*)]
\begin{multicols*}{3}
\item $sp$
\item $sp^2$
\item $sp^3$
\item $sp^3d$
\item $sp^3d^2$
\item $sp^3d^3$
\item $sp^3d^4$
\item $sp^3d^5$
\end{multicols*}\flushright  {\small Ans: (c)}
\end{enumerate}

\item  Indicate the hybridization of: \ce{CH4} 
\begin{enumerate}[label=(\alph*)]
\begin{multicols*}{3}
\item $sp$
\item $sp^2$
\item $sp^3$
\item $sp^3d$
\item $sp^3d^2$
\item $sp^3d^3$
\item $sp^3d^4$
\item $sp^3d^5$
\end{multicols*}\flushright  {\small Ans: (c)}
\end{enumerate}

\item  Indicate the hybridization of: \ce{H2O} 
\begin{enumerate}[label=(\alph*)]
\begin{multicols*}{3}
\item $sp$
\item $sp^2$
\item $sp^3$
\item $sp^3d$
\item $sp^3d^2$
\item $sp^3d^3$
\item $sp^3d^4$
\item $sp^3d^5$
\end{multicols*}\flushright  {\small Ans: (c)}
\end{enumerate}

\item  Indicate the hybridization of:  \hspace{.15in}\chemfig{\ce{C}(-[:-160]\ce{H})(-[:160]\ce{H})(=[:0]\ce{O})} 
\begin{enumerate}[label=(\alph*)]
\begin{multicols*}{3}
\item $sp$
\item $sp^2$
\item $sp^3$
\item $sp^3d$
\item $sp^3d^2$
\item $sp^3d^3$
\item $sp^3d^4$
\item $sp^3d^5$
\end{multicols*}\flushright  {\small Ans: (b)}
\end{enumerate}
\item  Indicate the hybridization of:  \hspace{.15in}\chemfig{\ce{C}(-[:180]\ce{H})(~[:0]\ce{C}(-[:0]\ce{H}))} 
\begin{enumerate}[label=(\alph*)]
\begin{multicols*}{3}
\item $sp$
\item $sp^2$
\item $sp^3$
\item $sp^3d$
\item $sp^3d^2$
\item $sp^3d^3$
\item $sp^3d^4$
\item $sp^3d^5$
\end{multicols*}\flushright  {\small Ans: (a)}
\end{enumerate}

\item  Indicate the hybridization of: \ce{XeF2} 
\begin{enumerate}[label=(\alph*)]
\begin{multicols*}{3}
\item $sp$
\item $sp^2$
\item $sp^3$
\item $sp^3d$
\item $sp^3d^2$
\item $sp^3d^3$
\item $sp^3d^4$
\item $sp^3d^5$
\end{multicols*}\flushright  {\small Ans: (d)}
\end{enumerate}


{\raggedright\textsc{\textbf{Molecular Orbital Theory}}\par}
\item  Using the MO order provided below, obtain the MO configuration for \ce{B2}: \\
$\sigma^{}_{2s}\sigma^{*}_{2s}\pi^{}_{2p}\sigma^{}_{2p}\pi^{*}_{2p}\sigma^{*}_{2p}$
 \begin{flushright}\small Ans:  $\sigma^{2}_{2s}\sigma^{*2}_{2s}\pi^{2}_{2p}$ \end{flushright}
\item  Using the MO order provided below, obtain the MO configuration for \ce{C2}: \\
$\sigma^{}_{2s}\sigma^{*}_{2s}\pi^{}_{2p}\sigma^{}_{2p}\pi^{*}_{2p}\sigma^{*}_{2p}$
 \begin{flushright}\small Ans:  $\sigma^{2}_{2s}\sigma^{*2}_{2s}\pi^{4}_{2p}$ \end{flushright}
 
\item  Using the MO order provided below, obtain the MO configuration for \ce{O2}:\\
$\sigma^{}_{2s}\sigma^{*}_{2s} \sigma^{}_{2p}\pi^{}_{2p}\pi^{*}_{2p}\sigma^{*}_{2p}$
 \begin{flushright}\small Ans:  $\sigma^{2}_{2s}\sigma^{2*}_{2s} \sigma^{2}_{2p}\pi^{4}_{2p}\pi^{2*}_{2p}$ \end{flushright}

\item  Using the MO order provided below, obtain the MO configuration for \ce{F2^+}: \\
$\sigma^{}_{2s}\sigma^{*}_{2s} \sigma^{}_{2p}\pi^{}_{2p}\pi^{*}_{2p}\sigma^{*}_{2p}$
 \begin{flushright}\small Ans:  $\sigma^{2}_{2s}\sigma^{2*}_{2s} \sigma^{2}_{2p}\pi^{4}_{2p}\pi^{3*}_{2p}$ \end{flushright}

%\item  Indicate the magnetic configuration of the molecule with MO configuration: $\sigma^{2}_{2s}\sigma^{2*}_{2s} \sigma^{2}_{2p}\pi^{4}_{2p}\pi^{3*}_{2p}$
%\begin{enumerate}[label=(\alph*)]
%\begin{multicols*}{1}
%\item paramagnetic
%\item diamagnetic
%\end{multicols*}
% \begin{flushright}\small Ans: (a) \end{flushright}



%
%\begin{center}\hspace{.05in}\chemfig{B(-[:0]\lewis{026,F})(-[:120]\lewis{024,F})(-[:240]\lewis{046,F})}\hspace{.05in}\hspace{.15in}\chemfig{\lewis{2,S}(=[:330]\lewis{06,O})(-[:210]\lewis{246,O})}\hspace{.05in} and \hspace{.05in}\chemfig{C(-[:0]H)(-[:90]H)(-[:180]H)(-[:270]H)}\hspace{.05in}
%\end{center}

\restoregeometry
\end{enumerate}
\end{multicols*}
\pagecolor{green!10}\afterpage{\nopagecolor}\newpage
\end{document}
